\documentclass[journal]{IEEEtran}

% --- Paquetes Requeridos y Recomendados ---
\usepackage[utf8]{inputenc}
\usepackage[T1]{fontenc}
\usepackage[spanish,es-tabla]{babel}
\usepackage{cite}
\usepackage{amsmath,amssymb,amsfonts}
\usepackage{graphicx}
\usepackage{textcomp}
\usepackage{xcolor}
\usepackage{booktabs}
\usepackage{microtype}
\usepackage{url}
\usepackage{hyperref}
\usepackage{tikz}
\usetikzlibrary{calc,babel}

% --- Paleta de Colores Oficial Universidad La Salle ---
\definecolor{lasalledark}{HTML}{0D2354}
\definecolor{lasalleblue}{HTML}{163982}
\definecolor{lasallelight}{HTML}{2E63C5}

% --- Gráficos e Hipervínculos ---
\graphicspath{{img/}}
\hypersetup{colorlinks=true,linkcolor=black,citecolor=blue,urlcolor=blue}

% --- Términos en español para IEEEtran ---
\addto\captionsspanish{\renewcommand\IEEEkeywordsname{Palabras clave}}
\renewcommand\IEEEkeywordsname{Palabras clave}

% ==============================================================================
% OBJETIVO DE EXTENSIÓN (requisito del curso)
% 5 a 10 páginas de SOLO TEXTO en IEEE dos columnas = ~4,250 a 8,500 palabras.
% Verificar con: make count   (o texcount -inc -sub=none main.tex)
% ==============================================================================

\begin{document}
\shorthandoff{<>}

% ==============================================================================
% CARÁTULA INSTITUCIONAL - UNIVERSIDAD LA SALLE (NO MODIFICAR)
% ==============================================================================
\begin{titlepage}
  \thispagestyle{empty}
  \begin{tikzpicture}[remember picture, overlay]
    \fill[lasalledark] (current page.north west) rectangle ([xshift=1.2cm]current page.south west);
    \fill[lasalleblue] ([xshift=1.2cm]current page.north west) rectangle ([xshift=1.5cm]current page.south west);
    \fill[lasallelight] ([xshift=1.5cm]current page.north west) rectangle ([xshift=1.7cm]current page.south west);
    \fill[lasalleblue!12] ([yshift=-4.5cm]current page.north east) -- (current page.north east) -- ([xshift=-6.5cm]current page.north east) -- cycle;
    \fill[lasalleblue!30] ([yshift=-2.8cm]current page.north east) -- (current page.north east) -- ([xshift=-4.0cm]current page.north east) -- cycle;
    \fill[lasallelight!50] ([yshift=-1.5cm]current page.north east) -- (current page.north east) -- ([xshift=-2.2cm]current page.north east) -- cycle;
    \fill[lasalleblue!15] ([xshift=-3cm]current page.south east) -- (current page.south east) -- ([yshift=2cm]current page.south east) -- cycle;
  \end{tikzpicture}
  \centering
  \begin{minipage}[c][0.92\textheight][s]{0.82\textwidth}
    \centering
    \vspace*{0.4cm}
    \includegraphics[width=0.38\textwidth]{img/logoLaSalleUniversidad.png}\par
    \vspace{0.4cm}
    {\large\bfseries\color{lasalledark} Facultad de Ingeniería\par}
    \vspace{0.15cm}
    {\large\bfseries\color{lasalleblue} Carrera Profesional de Ingeniería de Software\par}
    \vfill
    {\color{lasalleblue}\rule{\textwidth}{1.8pt}}\\[0.3cm]
    {\Large\bfseries\color{lasalledark} [TÍTULO PENDIENTE]\par}
    \vspace{0.3cm}
    {\large\bfseries\color{lasalleblue} Trabajo Final de Sistemas Operativos\par}
    \vspace{0.3cm}
    {\color{lasalleblue}\rule{\textwidth}{1.8pt}}
    \vfill
    \begin{flushleft}
      \hspace*{0.3cm}
      \begin{tabular}{ll}
        \textbf{\color{lasalledark} Asignatura:} & Sistemas Operativos \\[0.15cm]
        \textbf{\color{lasalledark} Docente:}    & Mg. Maribel Molina Barriga \\[0.15cm]
        \textbf{\color{lasalledark} Integrantes:} & $\bullet$ Diego Eduardo Zapana Choquehuayta \\
                                                  & $\bullet$ Heydy Anyely Castillo Arocutipa \\
                                                  & $\bullet$ Angelica Giovanna Huaman Cuba \\
                                                  & $\bullet$ Grecia Anhi Mamani Larico \\
      \end{tabular}
    \end{flushleft}
    \vfill
    {\large\bfseries\color{lasalledark} Arequipa, Perú\par}
    \vspace{0.15cm}
    {\small\color{gray} \today\par}
  \end{minipage}
\end{titlepage}

\clearpage
\twocolumn
\setcounter{page}{1}

% ==============================================================================
% ARTÍCULO CIENTÍFICO IEEE (DOS COLUMNAS) - ESQUELETO VACÍO
% ==============================================================================

\title{[TÍTULO PENDIENTE]\\[0.2cm]{\large Trabajo Final de Sistemas Operativos}}

\maketitle

% =========================================================================
% RESUMEN Y PALABRAS CLAVE
% =========================================================================
\begin{abstract}
% PENDIENTE (150--250 palabras, un solo párrafo, sin citas):
%   1) Contexto y motivación.
%   2) Problema / brecha de investigación.
%   3) Método o enfoque.
%   4) Hallazgos o resultados principales.
%   5) Conclusión e impacto.
\end{abstract}

\begin{IEEEkeywords}
% PENDIENTE: 4--6 términos normalizados (IEEE Thesaurus / ACM CCS).
\end{IEEEkeywords}

% =========================================================================
% I. INTRODUCCIÓN
% =========================================================================
\section{Introducción}
\label{sec:introduccion}

% PENDIENTE:
%   - Contextualización del área de Sistemas Operativos.
%   - Problemática técnica contemporánea.
%   - Brecha de investigación y motivación.
%   - Objetivos y preguntas de investigación (RQ1, RQ2, ...).
%   - Contribuciones.
%   - Organización del resto del artículo.

% =========================================================================
% II. MÉTODO
% =========================================================================
\section{Método}
\label{sec:metodo}

% PENDIENTE:
%   - Preguntas de investigación formalizadas.
%   - Estrategia de búsqueda y selección de literatura (2021--2026).
%   - Criterios de inclusión y exclusión.
%   - Taxonomía o criterios de comparación técnica.
%   - Proceso de análisis y síntesis.

% =========================================================================
% III. DESARROLLO Y DISCUSIÓN
% =========================================================================
\section{Desarrollo y Discusión}
\label{sec:desarrollo}

% PENDIENTE (sección central y más extensa):
%   - Subsecciones temáticas según la taxonomía.
%   - Análisis comparativo de mecanismos internos del kernel.
%   - Tablas y figuras de arquitectura.
%   - Discusión crítica: trade-offs, viabilidad, vacíos en la literatura.

% =========================================================================
% IV. CONCLUSIONES
% =========================================================================
\section{Conclusiones}
\label{sec:conclusiones}

% PENDIENTE:
%   - Respuesta directa a cada pregunta de investigación.
%   - Implicaciones prácticas.
%   - Limitaciones del estudio.
%   - Líneas de investigación futura.

% =========================================================================
% REFERENCIAS BIBLIOGRÁFICAS
% Se activan cuando exista la primera cita real en references.bib.
% =========================================================================
% \bibliographystyle{IEEEtran}
% \bibliography{references}

\end{document}
